\section{结果与讨论}

本章先报告 48 次独立重复运行与 200 个全题库消融沙箱的总体表现，再从研究轨迹中抽取问题，通过局部配对对照与内部状态干预分析性能差异的来源。

\subsection{端到端基准评估}

端到端评估由四题重复实验和全题库四配置实验组成，分别考察配置差异的重复性与跨任务覆盖。下文先报告两组实验，再以具体研究轨迹解释差异。

\subsubsection{四题三次重复的 48 次独立运行}

重复实验比较四种配置在 E1、NB1、HA1 和 IA1 上各三次独立运行，共 48 次。四题涉及消除预组织、多步脱羧、开环立体结果和方法边界。运行采用 DeepSeek-v4-flash，已归档对话、工具调用、计算账本与原生量化产物。重复运行检验配置差异是否依赖单次探索偶然性，并保留假设至交付的轨迹。

\discussionSpace

\subsubsection{全部 50 题的 200 沙箱四配置消融}

全题库实验将 50 题分别置于四种配置，每个任务—配置组合采用一个独立沙箱，共 200 个。它考察配置差异能否扩展到不同化学类别与研究要求，并与重复实验提供的稳定性信息相互补充。

\discussionSpace

\subsubsection{总体趋势与机制分析}

总体结果呈现出配置之间的差异：Guided 表现最强，Standard 的弱项主要集中在 Impossible 等任务；在 Easy 任务中，除 Standard 外，其余配置表现接近。下面结合研究轨迹，进一步分析这些差异与几何操作、证据组织及研究结束判断之间的关系。

\discussionSpace

下文从已归档轨迹抽取具体问题，以 Qwen3.8:27b 构造局部配对诊断，改变分子表示、选项顺序、几何信息、证据引用或推理模式。轨迹定位问题如何出现，配对输入检查哪些信息影响判断，内部状态干预追踪测量事实的作用，据此分析化学意图到空间操作、局部计算到科学证据的转换。

\subsection{从化学意图到可靠空间操作}

\subsubsection{化学身份为操作提供持续参照}
\begin{figure*}[tbp]
\centering
\includegraphics[width=\textwidth]{figures/fig06.pdf}
\caption{分子表示与选项顺序下的化学身份判断。16 对分子各包含四种输入条件；圆形和菱形分别表示正确与错误，内部字母保留模型实际回答。浅色区域为相同字符串条件。右侧显示第 14 对的两种 C\textsubscript{5}H\textsubscript{8} 连接结构及选项含义。}
\label{fig:06}
\end{figure*}


一次几何修改首先需要确定它作用于哪个分子、哪个原子和哪个反应中心。相同分子可以有不同的 SMILES 遍历方式；相同分子式也可以对应不同连接关系。为考察模型如何处理这些差别，我们构造了 16 对分子身份问题，覆盖相同身份、同分子式下的连接变化、连接与电荷位置变化、双键立体化学以及质子化变化。每一对交叉设置原字符串或等价改写，以及两种答案顺序，得到 64 次判断。

模型在其中 58 次作出正确判断。等价改写改变了 32 组表示配对中的 4 组语义判断，交换答案顺序改变了 32 组选项配对中的 6 组判断。后六组都重复选择字母 A。图~\ref{fig:06} 中的一对 C\textsubscript{5}H\textsubscript{8} 结构尤其直观：环内双键与环外双键对应不同连接方式，模型在四种条件下均回答 A，而 A 在两种顺序中分别代表“相同”与“不同”。这一结果将表示敏感性定位到实际的化学对象判断上，也说明持续研究需要保留稳定的身份参照。

Cloudra 的分子拓扑、原子映射与反应中心承担了这项职责。它们使模型在不同表述和不同构型间切换时，后续测量、片段移动及端点比较仍能指向原来声明的对象。原子对应关系还将产物端点对齐、路径插值与反应中心监测连接起来，使局部几何操作能够沿同一套化学身份继续进行。



\subsubsection{几何测量分担反复出现的数值推理}

识别研究对象之后，模型还需从坐标获得与操作有关的空间关系。我们选取 NB1 轨迹中的四个实际结构，要求比较两个指定原子间距。每个结构均交换答案顺序，并分别提供完整 XYZ 或程序测得的距离。两种信息条件进一步交叉直接作答与开启思考，形成 32 次原生聊天调用（图~\ref{fig:07}）。

直接作答时，XYZ 条件的八次响应全部选择 A，其中四次正确；提供测量距离后，八次均正确，并在选项交换后保持相同的语义答案。开启思考后，XYZ 条件得到七个正确答案，另一次在 1,024 个生成 token 时达到输出上限；该输入单独增加预算后，在 1,353 个生成 token 内得到正确答案。测量距离条件的八次回答均正确。开启思考的八组配对中，XYZ 与测量距离条件的生成 token 中位数分别为 807.5 和 86。

这些结果提示，几何判断的负担包含原子定位、数值运算和关系比较。增加推理可以完成这些局部计算，程序测量则直接提供后续判断所需的关系量。Cloudra 因而让模型以键长、键角、二面角与片段关系表达目标，由确定性算子执行坐标变化。这种分工将重复的数值工作交给程序，并使操作后的变化能够直接检查，与程序辅助语言模型的职责划分相呼应。\citep{ref15} 分支拥挤描述、反应物预组织与局部片段旋转在此基础上，进一步把空间关系用于候选构型的选择和构建。


\begin{figure*}[tbp]
\centering
\includegraphics[width=.64\textwidth]{figures/fig07.pdf}
\caption{坐标输入与测量距离输入下的几何判断。四个源结构分别在直接回答和开启思考条件下比较，每种条件均包含原始（F）与交换（R）选项顺序。圆形表示正确，菱形表示错误，三角形表示达到输出预算；内部字母保留实际回答。}
\label{fig:07}
\end{figure*}



\begin{figure}[tbp]
\centering
\includegraphics[width=\linewidth]{figures/fig08.pdf}
\caption{开启思考时的生成 token 数。每条线连接同一源结构、同一选项顺序下的笛卡尔坐标输入与测量距离输入，共八组配对比较。虚线表示输出上限。}
\label{fig:08}
\end{figure}


\subsubsection{反馈必须对应指定原子与几何目标}
\begin{figure*}[tbp]
\centering
\includegraphics[width=.9\textwidth]{figures/fig09.pdf}
\caption{E1 中面向指定目标的几何操作。Newman 投影保留原候选坐标沿 C2\(\to\)C3 观察时的方向。H7 为指定的转移原子；候选 A、B 的目标角分别为 59.4° 和 179.4°。投影旁列出相应操作顺序。}
\label{fig:09}
\end{figure*}


确定性执行还需要与化学意图保持对应。E1 的一条真实轨迹提供了完整例子：模型希望将转移氢 H7、C2、C3 与离去原子 Cl4 调整为接近 180° 的二面角。它围绕 C2–C3 轴分别旋转 C3 一侧片段 +120° 与 \( - \)120°，构建候选 A 和 B，并将 C3–Cl4 距离调整到 2.35 Å。默认结构报告中的代表性二面角涉及 H5，而当前目标涉及 H7。模型随后主动计算两候选的 H7–C2–C3–Cl4，获得约 59.4° 和 179.4°，并在下一次实际提交中采用候选 B。

图~\ref{fig:09} 给出由保存的候选坐标绘制的 C2\(\to\)C3 方向 Newman 投影，图~\ref{fig:10} 汇总对应的目标角残差。候选 B 的 H7 目标角距 180° 仅 0.6°，A 的相应残差为 120.6°；H5 比较量在两候选中的残差分别为 124.2° 和 115.8°。这组数据揭示了反馈的对象依赖性：一个读数的化学意义，来自它与当前指定原子及操作目的的对应关系。

围绕同一组测量值的局部对照进一步显示了信息与动作的连接。只提供 H5 时，三种循环选项排列均请求补测 H7；加入 H7 后，三种排列均选择 B。另两组充分性配对也随目标量的加入，从需要补测转为信息充分。由此，原子映射、内坐标、片段操作和反馈可以沿一条研究动作链理解：确定目标对象，执行空间改变，测量目标量，再选择下一候选。预组织、刚体对齐与路径准备也沿用这一对应关系，将反应意图传递到后续势能面探索。




\begin{figure*}[tbp]
\centering
\begin{minipage}[t]{.48\textwidth}
\vspace{0pt}\centering
\includegraphics[width=\linewidth]{figures/fig10.pdf}
\caption{候选 A、B 相对指定 180° 目标的圆周角残差。H7 为指定目标原子，H5 为比较原子。}
\label{fig:10}
\end{minipage}\hfill
\begin{minipage}[t]{.48\textwidth}
\vspace{0pt}\centering
\includegraphics[width=\linewidth]{figures/fig11.pdf}
\caption{三种局部选项排列下的后续动作。仅提供 H5 测量值时，模型请求 H7 测量；补充 H7 后，模型选择候选 B。}
\label{fig:11}
\end{minipage}
\end{figure*}




\subsubsection{测量事实的内部传递影响选择}

为进一步考察模型怎样使用已提供的测量值，我们在同一本地量化模型中进行了单位置激活替换。\citep{ref16} 输入中的一个距离由 2.029 Å 改为 1.029 Å，另一距离保持 1.248 Å，对应的比较答案随之反转。随后把原始输入在数字 token 位置的残差状态，分别替换到改动输入的第 7、15、31 和 47 层，并读取原始事实所对应答案的全词表概率。

替换的影响随层和选项顺序变化（图~\ref{fig:13}）。在交换选项的输入中，原始事实答案 A 的概率由改动后的 0.074，在第 7 和 15 层替换后分别升至 0.956 和 0.718；第 31 和 47 层的值分别为 0.072 和 0.070。在原选项顺序中，答案 B 的概率由 0.009，在前两层替换后分别变为 0.029 和 0.015，后两层仍接近改动输入。自替换及未改变前缀位置的控制替换中，全词表 logit 差均为零。

这一内部干预将“提供测量值”推进到“测量信息怎样参与选择”的因果机制层面：在交换选项的输入中，浅中层的状态替换显著反转了模型预测，证明特定位置的数值表征直接参与了后续决策；而在原选项顺序中，补丁未能冲破极低概率，这与 3.2.1 节揭示的行为级位置偏置（Positional Bias）高度自洽——当目标答案位于非偏好的选项 B 时，模型深层注意力可能被表层的选项先验吸附。这一不对称性从底层表征机制上提示，单纯在自由文本提示词中注入孤立数值，在模型长程推理中可能被注意力偏置与语义衰减削弱；依托确定性的结构化算子将几何测量、化学实体与操作语义强行绑定，可以在模型内部建立稳定抵抗先验偏置的科学决策回路。Cloudra 的结构化测量因此不仅是几何计算的结果，更是模型开展后续机理判断的物理锚点。


\begin{figure*}[tbp]
\centering
\includegraphics[width=.9\textwidth]{figures/fig12.pdf}
\caption{对单个报告距离的干预。一次数字修改将 C2–C4 距离从 2.029 改为 1.029 Å，O3–H8 保持 1.248 Å。在指定层，将原始输入中 token 位置 53 的残差状态替换到修改后输入的运行中。token 位置从零计数。}
\label{fig:12}
\end{figure*}



\begin{figure*}[tbp]
\centering
\includegraphics[width=.9\textwidth]{figures/fig13.pdf}
\caption{状态替换后原始事实答案的概率。曲线显示原始顺序（答案 B）与交换顺序（答案 A）下，原始事实答案在完整词表中的概率质量，包含无前导空格与带前导空格的 token 写法。水平线为原始输入与修改后输入的基线。层号从零计数。}
\label{fig:13}
\end{figure*}



\begin{figure*}[tbp]
\centering
\includegraphics[width=.75\textwidth]{figures/fig14.pdf}
\caption{状态干预实验的对照替换。在 p53 进行同状态自替换，以及在 p0 替换未改变的前缀，两种选项顺序下均保持全部词表 logits 不变。表中报告 logit 差值绝对值的最大值。token 位置从零计数。}
\label{fig:14}
\end{figure*}


\subsection{从局部计算结果到有证据的科学交付}

\subsubsection{DAG 保存结果之间的科学支持关系}

几何对象进入计算之后，研究需要把优化结构、能量、频率与路径逐步组织成对同一反应步骤的支持。每份结果的局部有效性及其归属共同决定这条支持链：频率应对应所声明的鞍点，IRC 应从该鞍点出发，两侧末帧的优化与频率应对应相应端点。DAG 将这些关系保存在物种与基元步骤上，为不断增加的结果提供持续的研究索引。执行状态、数值收敛与科学支持分别记录，使后续研究能够据此等待计算、调整重算或继续补充证据。

我们从归档提取四条源证据链，分别构造保持关联与改变单一引用的版本，并交叉平铺表示、关联表示和两种答案顺序，得到 32 个输入。单一引用的变化分别涉及频率几何、IRC 起始几何、端点优化输入或目标物种；局部终止、收敛、能量和频率值保持不变。图~\ref{fig:15} 以 E1 为例：将频率的几何引用从鞍点改为一个已有的端点，便改变了这份频率对当前步骤的支持关系。

关联表示的 16 次判断均正确，平铺表示的 16 次中有 14 次正确；两处错误均出现在引用已改变且选项交换的条件，分别来自 E1 与 NB1 第二步。随后选择的两条源链原生跟进中，开启思考使全部 16 次关联判断正确，包括原先两处错误。这一结果提示，明确关联可以分担模型查找与拼接计算关系的工作，也与图编码方式影响语言模型关系推理的观察相呼应。\citep{ref17} 在 Cloudra 中，这种表示进一步与真实计算产物绑定，使同一物种、步骤和尝试下的科学事实能够持续累积。

\begin{figure*}[tbp]
\centering
\begin{minipage}[t]{.48\textwidth}
\vspace{0pt}\centering
\includegraphics[width=\linewidth]{figures/fig15.pdf}
\caption{E1 源链中的证据与几何关联。受控输入将频率计算的几何引用从 g\_63 改为 g\_42，同时保持局部终止、收敛、能量和频率数值。图中保留从过渡态优化、双向 IRC 到端点优化与频率计算的关联。}
\label{fig:15}
\end{minipage}\hfill
\begin{minipage}[t]{.48\textwidth}
\vspace{0pt}\centering
\includegraphics[width=\linewidth]{figures/fig16.pdf}
\caption{四条源链共 32 个输入的关联判断，交叉比较一致或修改后的关联、平铺或关联式表示，以及原始（F）或交换（S）答案顺序。字母保留已保存的 A/B 回答；圆形表示正确关联，菱形表示错误关联。}
\label{fig:16}
\end{minipage}
\end{figure*}

\begin{figure*}[tbp]
\centering
\includegraphics[width=.96\textwidth]{figures/fig17.pdf}
\caption{原始 IA1 轨迹中的完成请求与 DAG 反馈。报告仍有未解决的证据后，模型请求完成；系统指出缺少 short\_irc 证明并向模型反馈，随后模型改为请求部分完成。}
\label{fig:17}
\end{figure*}

\subsubsection{诊断将缺口转化为具体研究动作}

证据链的断裂状态需要转化为具有明确物理针对性的反馈，才能有效驱动模型采取纠偏动作。若仅向模型暴露未过滤的量子化学原生输出，冗长的数值日志极易导致大语言模型在海量上下文中忽略关键的几何与物理异常；而结构化诊断算子通过消除冗余信息，将底层失败精准映射为具体的物理缺陷（如空间位阻冲突原子对、虚频非目标对齐或路径端点脱靶），为下一轮局部尝试提供了明确的操作靶标。

真实探索轨迹清晰展现了诊断颗粒度对自愈行为的重要影响。在 Brief-errors 的 HA1 案例中，系统按消融设计移除了关于 IRC 端点连接缺陷的具体诊断，仅保留宽泛的报错代码，导致计算智能体陷入盲目重试，仅仅机械更换源极小值引用而始终未能修复内在的几何连接缺口；与之形成鲜明对照的是，在 Guided 的 NB1 案例中，底层检查捕捉到零虚频状态并给出模式错误分类，协调智能体据此迅速拟定针对性的重试策略，引导计算智能体调整反应初猜与鞍点优化约束，并将改进精准落实到新提交的计算作业中。这一案例对比可以证明，长程机理探索中的自愈恢复不仅取决于“是否知道计算受挫”，更取决于系统能否将物理证据缺陷精准桥接至“应当修改哪个化学对象或计算参数”。\citep{ref18}

这一机制完整诠释了 Cloudra 多智能体协同的实证优势：计算智能体依托局部的几何与路径反馈迭代具体计算尝试，协调智能体统揽分支全局并权衡是否切换机理假说，督导智能体则对照法定契约独立把关。由中央调度账本支撑的异步执行与证据持久化，支持各角色在长周期任务中持续围绕同一化学实体稳定推进，减少了对话历史膨胀带来的信息漂移。

\subsubsection{未完成项必须持续参与结束判断}

机理交付的最终关卡在于全局判断：当前已积累的计算产物是否真正满足了严苛的物理化学验收契约。IA1 的真实自主探索轨迹直观揭示了大语言模型在复杂科研推理中的固有缺陷——早熟完工幻觉（Premature Completion Hallucination）。在该案例中，关键基元步尚缺少 short\_irc 双向路径确证，且豁免清单为空；然而，由于督导智能体在过程对话中曾对部分驻点优化的收敛给出了阶段性肯定，模型便产生认知漂移，试图据此直接宣称整个反应机理研究已经“完全完成”（complete）。若缺乏独立的底层硬门禁，此类在物理化学上缺乏充分证据支撑的机理主张将轻易逃逸并造成虚假交付。

在此处，交付图底座确立的确定性物理门禁发挥了决定性的拦截作用。系统直接驳回了模型的提前完工请求，并将缺少 short\_irc 证据链的物理缺口强制注回推理上下文；面对无法逾越的刚性事实，模型立即在随后的提交中校正立场，将结束状态修正为符合物理现状的“部分完成”（partial，图~\ref{fig:17}）。

当物理状态被显式固化为不可篡改的结构化账本时，大模型便能够稳定遵循这一事实源开展推理。在围绕 IA1 构造的 12 次上下文窗口测试（按正反选项顺序组成 6 对）中，无论是仅提供最新结构化报告，还是引入包含 12 条消息的历史窗口，模型在所有状态提问与选项逆转条件下均百分之百准确复核出尚未满足的科学缺口。交付 DAG 因而构成了模型后续推演不可动摇的物理坐标系：明确界定哪些对象已获支撑、当前缺口位于何处，以及何时应当追加计算或据实交付阶段性成果。

\subsection{工作流引导、局限与展望}

Guided 在基准测试中取得了最强的总体表现，但在 Easy 任务中并未与其余非 Standard 配置明显拉开差距。当任务需要持续探索反应路径、组织计算证据并从失败计算中恢复时，其优势更为明显。这一分布说明，Cloudra 工作流引导的价值与研究任务的要求有关：基础任务往往可以通过多种路径完成，而更复杂的任务更依赖几何操作、计算证据和结束判断在探索过程中保持衔接。

Record 提供了一个有启发的对照。它保留了 Guided 可调用的工具与诊断能力，但移除了强制科学验收；其在 Normal 和 Hard 任务中的表现仍落后于 Guided。在观察到的轨迹中，模型并未始终主动调用这些能力或遵循推荐工作流，而是选择了临时拼接的替代路径。因此，工具可用并不意味着模型会主动利用这些工具构建完整的科学论证。这一对比支持将工具能力与明确的研究要求相结合，并通过反馈持续引导模型处理尚未解决的证据缺口。

Brief-errors 呈现出类似的启示。压缩诊断后，部分轨迹更容易反复提交计算而未解决底层缺陷，上述 HA1 端点关联案例即为一例。由于缺少足够具体的引导将研究拉回预期工作流，重复尝试可能持续消耗算力，而科学缺口依然存在。其低于 Guided 的表现与具体诊断的重要性一致：反馈需要同时指出哪一项要求未满足，以及应采取什么研究动作来解决。Record 与 Brief-errors 的表现共同支持了 Cloudra 的系统设计理念，即将化学意图连接至确定性操作，并以持续保留的证据要求引导后续研究。

Standard 的观测运行时间最短，但这种速度未转化为相应的基准得分。Standard 在正式基准运行中频繁遗漏交付图，交付不完整是该配置一个突出且反复出现的问题。我们在离线整理中，仅利用对应原运行已有的研究记录与计算产物补建缺失的交付图，没有增加计算或新的科学结论。补建交付图后，其得分仍然较低，说明性能差距不能仅由交付格式缺失解释。这一对比也说明，运行时间需要与实际交付的、获得科学证据支持的工作量共同理解。

本基准存在与计算资源有关的局限。由于可用算力有限，我们采用了较高的沙箱并发，CPU 资源排队可能延迟个别计算，减少一次运行在既定预算内完成的计算与证据量，进而影响得分。因此，当前结果同时反映了研究行为与共享计算环境的影响；CPU 排队是得分差异的一个可能因素，尚不能据此解释某一配置的具体失分。后续评估中进一步区分排队时间与实际计算时间，有助于明确这一因素的影响。

Cloudra 缓解了语言模型直接进行几何推理与操作的困难，其构象采样能力相较早期版本也已有明显提升，但高柔性体系仍然具有挑战。随着构象多样性增加，寻找相关的反应物排布、中间体和反应路径可能需要更多采样与量子化学计算。当前系统尚不足以在较小算力预算内，独立处理复杂程度显著提高的高柔性体系。提高构象探索效率，并在相互竞争的候选结构之间更合理地分配计算资源，是后续发展的重要方向。

我们的长期目标是将 Cloudra 从小分子机理计算扩展至更广泛的计算化学领域。当前版本对多种计算引擎的支持、共享执行基础设施及持续保存的科学记录，已为这一扩展提供了基础。面向新的问题类型，还需要配套的科学表示、计算方法与验证准则。我们希望 Cloudra 最终成为计算化学家的强大研究工具，以及实验团队的计算专家。Cloudra 将以 GNU Affero General Public License（AGPL）开源，支持系统的公开检查、复用与社区共同发展。
